% discussion.tex (v11)
\section{Discussion}\label{sec:discussion}

\subsection{Physical anchoring and transferability}
The results show that the way physical knowledge enters the model determines whether ML improves ALT prediction beyond a regional Stefan model and whether the improvement transfers to new regions. In the physics-residual formulation (equation~(\ref{eq:residual})), the Stefan term carries the response to the thaw-season climate of each year, which extrapolates to new regions and years in a physically bounded form, and the learner corrects departures of the effective coefficient associated with soil and terrain. Four results support this division of roles: the gain required terrain and soil covariates; it was largest where the Stefan error was largest; the attribution of the residual learner pointed to terrain in the Lena Delta and to soil organic carbon in Canada; and Stefan pseudo-labels pulled the learner towards the Stefan model without adding accuracy. Direct ML, which must learn the climate response from the labels, reached lower point errors than the physics-residual model when all labels of a region were available, but its Lena Delta error exceeded that of the Stefan model by \ndrefLeregDzabs{} cm when the region was withheld. This contrast agrees with studies of physics-guided learning in limnology and hydrology, in which process knowledge mainly improves performance under data scarcity and extrapolation\cite{Read2019,Jia2021,Willard2022}. Because most permafrost regions have few or no ALT observations, the accuracy that a method retains without target-region labels is more relevant for mapping than its accuracy with all labels of a well-observed region.

\subsection{Comparison with previous studies}
Previous ALT mapping studies differ from the present one in the climate input, the validation design and the reference used for comparison. ML products relate observations pooled over many years to multi-year predictors and report random or site-based cross-validation on their own site sets\cite{Ran2022a,Liu2024,Wei2026}. In our data, random cross-validation underestimated the spatial-block error of direct ML by \noptimLeDz--\noptimCaDz{} cm, and withholding the Lena Delta increased its error there to \nrmseLeregDz{} cm, compared with \nrmseLeregPz{} cm for the Stefan model with the coefficient of the other regions. Errors from random or site-based splits therefore do not indicate the accuracy of an ML product in a region without training observations. Gautam et al.\cite{Gautam2025} compared a random forest with the Stefan model at 68 CALM sites in Alaska observed in 2014; the coefficient of determination of the random forest fell from 0.84 in training to 0.24 on a random test split, whereas the Stefan model retained 0.54. The present study extends this comparison to site-years of 2010--2024 in three regions and to region hold-out, and combines the two approaches so that the learner corrects the Stefan model instead of replacing it. Process-based products such as ESA CCI Permafrost are not trained on ALT observations\cite{Westermann2024,Westermann2025}; at the external CALM sites, the error of ESA CCI v5 (\nftAllCcifivezero{} cm) exceeded that of our model with five local labels (\nftAllRsfive{} cm).

Scoring all models and products on the same cells and years, with masks for product training sites, makes these differences comparable. The regional ranking of the products followed their training data. Where the products were trained on nearby observations (Alaska and Canada), they were as accurate as our model; where they were not (the Lena Delta), they overestimated thin active layers and were less accurate than a model trained without any Lena Delta label. The residual learner also improved a product used as the anchor (equation~(\ref{eq:prod})), so that local labels combined with residual learning can adapt both a physical model and an existing product to a region that the product did not learn. The same-cell, same-year protocol and the training-site masks can be applied to other products and regions.

\subsection{Workflow for new regions and years}
The label-number experiments translate into a sequence for applying the method to a new region. Without local labels, the Stefan model with the coefficient of other regions, or the physics-residual model trained on other regions, is a safer starting point than direct ML, whose Lena Delta error without target labels was \ntlrmseLeDszero{} cm in the transfer experiment (Fig.~\ref{fig:labels}b). Twenty-five Lena Delta labels reduced the error of the physics-residual model from \ntlrmseLeRszero{} to \ntlrmseLeRstwfive{} cm, below that of the source-coefficient Stefan model (\ntlrmseLePzzero{} cm), and in the external CALM regions five labels reduced the error from \nftAllRszero{} to \nftAllRsfive{} cm. These results agree with the value of a limited number of local observations reported for regionalisation in hydrology\cite{Pool2019}. The number of labels required depends on the distance of the new region from the training range: in Canada, where thaw depths vary widely, recalibration on 25 labels increased the error, and 25--100 labels that span the local range of soils and terrain are a reasonable target. Because each map is driven by the ERA5-Land climate of a stated year (TDD to 1 September) and carries an area of applicability and a prediction interval, the same workflow produces yearly maps for monitoring, allows comparison with observations or products of the same year, and marks the areas outside the area of applicability, where additional observations are needed.

\subsection{Limitations and outlook}
Half of the Alaska labels come from 2014, which prevented a resolved estimate of the effect of the year-specific input in that region. The Lena Delta labels come from one data set in \nnblkLe{} blocks, which widens the confidence intervals. The climate input is monthly and the soil covariates are static, so that short-term weather and changes in soil moisture are not represented. Active layers thicker than about 3 m, mainly in mountains and at the southern margin of the permafrost zone, lie outside the training range, and the external tests show that the residual learner does not transfer to such conditions without local labels. Daily climate inputs, information on snow and soil moisture, and labels from a fourth region such as West Siberia, where CALM and campaign data both exist, are the next steps.
